\chapter{Caracterizacion termica}
\section{Prueba del maximo seguro}
Para caracterizar termicamente al sistema se necesita conocer la curva de enfriamiento que se estabiliza cerca del punto de operacion deseado.
La curva se encuenrta al aplicar pequeños escalones seguros de corriente $\Delta I_{0}$ con el objetivo de enfriar, la entrada electrica y la respuesta termica del sistema deben registrarse a lo largo del tiempo.
Esta curva nos permite conocer la  ganancia termica K, la constante de tiempo termica efectiva $\tau$ y el retardo L del sistema, sustituir esta informacion en la ecuacion (2.44) permite conocer el modelo a primer orden de la planta termica.
Mediante las ecuaciones (2.79), (2.82) y (2.94) se pueden encontrar las constantes $K_{p}$ y $K_{I}$ para diseñar el control proporcional del sistema.
\section{Caracterización térmica del C30902SH-2 DTC}